%%%%

%%%   Самарский государственный технический университет

%%%   Кафедра прикладной математики и информатики

%%%

%%%   Преамбула для статьи в журнал "Вестник Самарского государственного

%%%   технического университета. Серия: «Физико-математические науки»"

%%%   Ориентирована под MikTEx 2.4 for Win32

%%%

%%%   Хотя сам журнал собирается под GNU/Linux на дистрибутиве TeXLive



\documentclass[11pt,a4paper,twoside]{article}



\usepackage[cp1251]{inputenc}

\usepackage[T2A]{fontenc}

\usepackage[english,russian]{babel}

\usepackage{samgtubib}

\usepackage[dvips]{graphicx}

\usepackage[multiple]{footmisc}



\usepackage{samgtu}

\renewcommand{\VestnikYear}{2009} % Год издания Вестника

\renewcommand{\VestnikNumber}{2\,(19)} % Номер Вестника

\renewcommand{\VestnikMonth}{Ноябрь} % Месяц выхода Вестника

\renewcommand{\VestnikData}{01.11.09} % Дата подписания Вестника в печать

\renewcommand{\VestnikUPL}{26,64} % Усл. печ. л.

\renewcommand{\VestnikUIL}{25,73} % Уч.-изд. л.

\renewcommand{\VestnikRegNumber}{479/09} % Регистрационный номер

\renewcommand{\VestnikOrder}{994} % Номер заказа






%
%\usepackage{pst-barcode}
%\usepackage{auto-pst-pdf}
%
%%
%%\newcommand{\totop}[1]{\raisebox{6pt}[1pt][2pt]{#1}} 
%%\newcommand{\tottop}[1]{\raisebox{12pt}[1pt][2pt]{#1}} 
%%\newcommand{\totttop}[1]{\raisebox{18pt}[1pt][2pt]{#1}} 
%%\newcommand{\tottttop}[1]{\raisebox{24pt}[1pt][2pt]{#1}} 
%%\newcommand{\totttttop}[1]{\raisebox{30pt}[1pt][2pt]{#1}} 
%%\newcommand{\tobot}[1]{\raisebox{-6pt}[1pt][2pt]{#1}} 
%%\newcommand{\tobbot}[1]{\raisebox{-12pt}[1pt][2pt]{#1}} 
%%\newcommand{\tobbbot}[1]{\raisebox{-18pt}[1pt][2pt]{#1}}
%
%
%%\samgtupubl % для генерации pdf, отдаваемого в типографию СамГТУ
%\pdfcrop % обрезка pdf для размещения на math-net.ru
%\draft % На корректуре
%\filllastpage % Последняя страница не заполнена не полностью
%\briefcomm % Статья является кратким сообщением

\vsgtu{1352}

\FirstPage{170}
\LastPage{177}
%
%
%\begin{document}




%\hfill \begin{pspicture}(1in,1in)
%\psbarcode{http://dx.doi.org/10.14498/vsgtu1352}{}{qrcode}
%\end{pspicture}
%\vspace{-1in}
%
%\Partition{Теоретическая физика}{Theoretical Physics}

\renewcommand{\baselinestretch}{0.95}


\udk{53.03:(539.183--539.194)}
\msc{81S40, 58D30}

\rutitle{Специфика перехода к мнимому времени в~интеграле по траекториям при описании коллапса волновой функции}
{Представление коллапса волновой функции в комплесном времени}

\ruauthor{Н.~В.~Мелешко, А.~Ю.~Самарин}
 {М\,е\,л\,е\,ш\,к\,о~Н.~В.\, С\,а\,м\,а\,р\,и\,н~А.~Ю.}

\enauthor{N.~V.~Meleshko, A.~Yu.~Samarin}{M\,e\,l\,e\,s\,h\,k\,o~N.~V.,  S\,a\,m\,a\,r\,i\,n~A.~Yu.} 
 
 

\ruaffil{\rusamgtu.}

%\email{1}{\mem{samarinay@yahoo.com}(A.Yu.~Samarin, \CAut)}


\ruauthordetails{2}{\fullauthor{\rupid{60526}{Наталья Витальевна Мелешко}}\regalia{\mem{meleshko1958@gmail.com}}\position{старший преподаватель} \dept{каф. общей физики и физики нефтегазового производства}\\
\fullauthor{\rupid{42489}{Aлексей Юрьевич Самарин}}\regalia{к.ф.-м.н, доц.; \mem{samarinay@yahoo.com}; автор, ведущий переписку}\position{доцент}\dept{каф. общей физики и физики нефтегазового производства}}

\enauthordetails{5}{\fullauthor{\enpid{60526}{Natalia V. Meleshko}} \regalia{\mem{meleshko1958@gmail.com}} \position{Senior Lecturer} \dept{Dept. of General Physics and Physics of Oil and Gas Production}
\\ 
\fullauthor{\enpid{42489}{Alexey Yu. Samarin}} \regalia{Cand. Phys. \& Math. Sci.; \mem{samarinay@yahoo.com}; Corresponding Author}\position{Associate Professor} 
\dept{Dept. of General Physics and Physics of Oil and Gas Production}}

\rukeywords{коллапс волновой функции, интеграл по траекториям, интегральная мера Винера, комплексное время, поворот Вика}

\ruabstract{С помощью представления времени в комплексной форме интеграл \linebreak по~траекториям, описывающий временное изменение волновой функции квантовой частицы, преобразуется к вещественному виду. Подобная процедура необходима для того, чтобы иметь возможность присвоить меру множествам виртуальных траекторий в континуальном интеграле, определяющем амплитуду перехода между квантовыми состояниями. Квантовая амплитуда перехода в форме континуального интеграла является вещественной функцией модуля комплексного времени для всех мнимых значений последнего. При этом отрицательные значения времени соответствуют обратной последовательности событий. В силу обратимости законов механики при описании квантовой эволюции в виде движения индивидуальных точек по виртуальным траекториям знак времени не имеет значения. Поэтому при определении интегральной меры рассматривается отрицательная полуось мнимого времени, для которой квантовый интеграл по траекториям имеет вид интеграла Винера с мерой в форме интеграла Эйнштейна"--~Смолуховского. Что касается описания коллапса волновой функции, то из-за его необратимости во времени последовательность событий, определяемая модулем комплексного времени, не должна нарушаться. Вследствие этого указанное преобразование интеграла для описания коллапса может производиться только в~верхней полуплоскости комплексного времени. Показана возможность аналитического продолжения меры Винера на верхнюю полуплоскость комплексного времени, что обеспечивает существование квантового континуального интеграла для любых значений комплексного времени. Это обстоятельство позволяет представить интеграл по траекториям в вещественном виде как функцию мнимого времени для прямой последовательности событий. Тем самым появляется возможность учесть влияние имеющей место при коллапсе скачка потенциальной энергии в функционалах действия на вес соответствующих виртуальных траекторий.}

\entitle{Complex time transformations peculiarities for wave function collapse description using quantum path integrals}



\enaffil{\ensamgtu.}


\enabstract{A quantum path integral was transformed into the real form using a complex representation of the time. Such procedure gives the possibility to specify measures for the sets of the virtual paths in continual integrals determining amplitudes of quantum states transitions. The transition amplitude is a real function of the complex time modulus. Negative time values correspond to the reverse sequence of events. The quantum evolution description in form of the virtual paths mechanical motion does not depend on the sign of the time, due to the reversibility of the classical mechanics laws. This allows to consider the negative half of the imaginary axis of the time for the path integral measure determination. In this case this integral has the form of Wiener's integral having the well-known measure. As the wave function collapse is irreversible effect, the causal chain of events cannot be changed. Thus, to describe the collapse the transformation of quantum path integrals have to be performed in upper half plane of the complex time. It is shown that the Wiener measure for the real continual integral can be continued analytically on this actual range of the complex time. This allows to use the quantum path integral for any actual range of the complex time.}

\enkeywords{wave function collapse, path integral, Wiener measure, complex time, Wick rotation}

\date{13/XI/2014} % Дата поступления статьи в редакцию.
\revisiondate{26/XI/2014} % В окончательном виде
\accepteddate{27/XI/2014}

\makerutitle


 \Section[n]{Введение} Существуют два способа изменения механического состояния квантового объекта, которые традиционная квантовая механика определяет как эволюцию и редукцию волновой функции \cite{bib:E}. В рамках этой теории эволюция и редукция постулируются независимо друг от друга и не имеют никакого общего физического основания. Более того, редукция вообще не рассматривается как физический процесс, обусловленный причинно"=следственными связями между событиями [\citen{bib:A}--\citen{bib:C}]. В работах~\cite{bib:D,bib:L} показана возможность описания обоих механизмов преобразования волновой функции на основе единого динамического закона. В качестве такого закона было использовано интегральное волновое уравнение с ядром в виде интеграла по траекториям. Отказ от постулата <<наблюдаемой>> при использовании этого уравнения для описания процессов временного изменения механического состояния квантовой частицы позволил вывести его из структуры так называемого фейнмановского представления квантовой механики \cite{bib:F,bib:KV} и предложить в качестве динамической основы более общей теории. Эта теория рассматривает оба способа изменения состояния квантовой системы как два типа проявления единого механизма преобразования волновой функции в различных условиях. Эволюция при этом может рассматривается как механическое движение сплошной среды квантовой частицы  \cite{bib:V}, а редукция "--- как нелокальное преобразование внутренней структуры самой квантовой частицы в процессе измерения ее характеристик \cite{bib:PRL}.
 В последней работе утверждается, что физической причиной пространственной локализации квантовой частицы при её взаимодействии с макроскопическим измерительным прибором является бесконечный локальный скачок потенциальной энергии квантовой частицы до макроскопического значения. Чтобы оценить последствия такого скачка, необходимо преобразовать квантовый интеграл по траекториям к вещественному виду [\citen{bib:G}--\citen{bib:K}]. Для этой цели время представляется в комплексном виде и интеграл по траекториям рассматривается для мнимой оси комплексного времени. Последняя процедура может быть формально осуществлена двумя способами. Так, при определении меры континуального интеграла \cite{bib:EDK} и описании механического движения центра масс макроскопического тела~\cite{bib:KSV} поворот в комплексной области осуществляется по часовой стрелке, а при описании процесса редукции "--- против  \cite{bib:PRL}. Причины указанного выбора и его правомерность при этом не указываются. В настоящей работе проводится анализ правомерности использования аналитического продолжения интеграла по траекториям на положительную и отрицательную полуоси мнимого времени как с физической, так и с математической точек зрения.

\smallskip

 \Section[n]{Направление вращения}
Рассмотрим одномерное движение квантовой частицы. Под термином <<квантовая частица>> понимается в общем случае распределенный в пространстве объект, который в процессе пространственной локализации преобразуется в материальную точку. В соответствии с \cite{bib:F,bib:KV} это движение описывается интегральным волновым уравнением вида
 \begin{equation}\label{eq:math:ex1}
    \Psi_{t_{2}}(x_{2})=\int K(x_{1},x_{2})\Psi_{t_{1}}(x_{1})\,dx_{1},
\end{equation}
где $ x_{1} $, $t_{1} $, $ x_{2} $, $t_{2} $ "--- координаты пространства и моменты времени, соответствующие начальной и конечной волновым функциям; $ \Psi_{t_{1}}(x_{1})$, $ \Psi_{t_{2}}(x_{2})$ "--- волновые функции начального и конечного состояний; $ K_{t_{1},t_{2}}(x_{1},x_{2})$ "--- амплитуда перехода, которая записывается в форме интеграла по траекториям:
\begin{equation}\label{eq:math:ex2}
    K_{t_{1},t_{2}}(x_{1},x_{2})=\int \exp{\frac{i}{\hbar}S_{1,2}[x(t)]}\,[dx(t)],
\end{equation}
где $ S_{12}[\gamma]$ "--- классический функционал действия для виртуальной траектории $ \gamma\equiv x(t)$, которая имеет начало в точке пространства $ x_{1}$ и конец в точке $ x_{2}$ в~моменты времени $ t_{1}$ и $ t_{2}$ соответственно; $[d\gamma]$ "--- элемент объема пространства виртуальных траекторий \cite{bib:EDK}. Чтобы определить меру множеств траекторий в последнем интеграле, его преобразуют к вещественному виду. Для этого время формально представляется в комплексном виде: $ t=\tau e^{i\varphi} $, где $\tau $ и $\varphi $ "--- вещественные величины (для реальной физической ситуаций $\varphi=0 $). Далее рассматривается поворот в комплексной плоскости времени на угол ${\pi}/{2} $ по часовой стрелке. При таком повороте для кинетической энергии имеем
\begin{equation}\label{eq:math:ex4}
  T=-\frac{m}{2}\Bigl(\frac{dx}{d\tau}\Bigr)^{2}.
\end{equation}
В случае отсутствия явной зависимости потенциальной энергии от времени амплитуда перехода примет следующий вид:
\begin{equation}\label{eq:math:ex5}
    K_{t_{1},t_{2}}(x_{1},x_{2})=\int \exp\Biggl({\frac{1}{\hbar}\int\biggl(-\frac{m}{2}\Bigl(\frac{dx}{d\tau}\Bigr)^{2}}-U(x)\biggr)d\tau\Biggr)[dx(\tau)].
\end{equation}
Этот интеграл формально совпадает с интегралом Винера для броуновского движения, мера множеств траекторий которого определяется интегралом Эйнштейна"--~Смолуховского~\cite{bib:KPT}. Тогда значение квантового интеграла по траекториям может быть получено в вещественном виде как функция мнимого отрицательного времени. Искомая величина амплитуды перехода в этом случае определяется аналитическим продолжением интеграла Винера на вещественную ось времени. Математическая корректность этой процедуры не вызывает сомнений, однако переход к отрицательному знаку перед модулем времени меняет направление течения физических процессов. 
Для процесса эволюции волновой функции, в течение которого набор виртуальных траекторий частицы остается неизменным, это изменение сводится исключительно к перемене пределов в интегралах эвклидова действия (для мнимой оси времени) на каждой из виртуальных траекторий. То есть континуальный интеграл (амплитуда перехода) однозначно определен как функция времени как для прямого, так и для обратного течения событий. Это обстоятельство позволяет при определении меры квантового интеграла по траекториям пользоваться аналитическим продолжением амплитуды перехода на отрицательную часть мнимой оси времени. Целесообразность такой процедуры обусловлена совпадением вещественного квантового интеграла по траекториям с~интегралом Винера, для которого мера хорошо известна.

Иная ситуация имеет место при возникновении редукции волновой функции в процессе перехода квантовой частицы из одного состояния в другое. Согласно \cite{bib:PRL}, при коллапсе происходит изменение пространства виртуальных траекторий. Это изменение однозначно определено предысторией движения системы. После коллапса всякая информация об <<исчезнувших>> траекториях частицы теряется.
Таким образом, для обратного хода времени определение временной зависимости амплитуды перехода при наличии коллапса принципиально невозможно. Это обстоятельство делает необходимым при анализе процесса редукции рассматривать только верхнюю полуплоскость комплексного времени в качестве области определения амплитуды перехода.

\smallskip

\Section[n]{Аналитичность континуального интеграла в комплексной плоскости}
 Область определения интеграла по траекториям $0<\varphi<\pi$ ранее не анализировалась с точки зрения его аналитичности. Интегральная мера определяется для $\varphi={-\pi/2}$. Элементарные множества виртуальных траекторий (квазиинтервалы), для которых задается мера, определяются условиями
 \begin{equation}\label{eq:math:ex5}
   \bigl\{\alpha_{1}<x(\tau_{1})<\beta_{1},\alpha_{2}<x(\tau_{2})<\beta_{2},\dots,\alpha_{n}<x(\tau_{1})<\beta_{n}\bigr\}.
\end{equation}
Квазиинтервалам присваивается мера Винера, определяемая интегралом Эйнштейна "--- Смолуховского~\cite{bib:KPT}
\begin{multline}\label{eq:math:ex6}
\frac{1}{\bigl(\sqrt{2\pi\tau}\bigr)^{n}}\int_{\alpha_{1}}^{\beta_{1}}\dots
\int_{\alpha_{n}}^{\beta_{n}}
\exp\Bigl(-\frac{(x_{1}-x_{0})^{2}}{\tau_{1}-\tau_{0}}\Bigr)
\exp\Bigl(-\frac{(x_{2}-x_{1})^{2}}{\tau_{2}-\tau_{1}}\Bigr)
\cdots\\
\cdots\exp\Bigl(-\frac{(x_{n}-x_{n-1})^{2}}{\tau_{n}-\tau_{n-1}}\Bigr)dx_{1}\cdots\,dx_{n-1}.
\end{multline}
Отличие квантового интеграла по траекториям от интеграла Винера состоит в том, что для интеграла Винера время $t=\tau$, тогда как для квантового интеграла $t=-i\tau$. Для перехода к реальному времени необходимо аналитически продолжить интеграл по траекториям на вещественную временную ось. Для $-\pi<\varphi<0 $ квантовый интеграл по траекториям имеет следующий вид:
\begin{multline*}
    K_{t_{1},t_{2}}(x_{1},x_{2})=\int \exp\Biggl(\frac{i}{\hbar}\int\biggl(\frac{m}{2}\Bigl(\frac{dx}{d\tau}\Bigr)^{2}-U(x)\biggr)\cos\varphi d\tau+
    \\
    +\frac{1}{\hbar}\int\biggl(\frac{m}{2}\Bigl(\frac{dx}{d\tau}\Bigr)^{2}-U(x)\biggr)\sin\varphi d\tau\Biggr)[dx(t)].
\end{multline*}
В этой области $\sin\varphi$ отрицателен. Отрицательный знак перед кинетической энергией частицы, выраженной через скорость ее движения по виртуальной траектории, во втором слагаемом показателя экспоненты обеспечивает аналитичность функции
\begin{multline}\label{eq:math:ex7}
\frac{1}{\bigl(\sqrt{2\pi\tau}\bigr)^{n}}\int_{\alpha_{1}}^{\beta_{1}}\dots
\int_{\alpha_{n}}^{\beta_{n}}
\exp\Bigl(-\frac{(x_{1}-x_{0})^{2}}{t_{1}-t_{0}}\Bigr)
\exp\Bigl(-\frac{(x_{2}-x_{1})^{2}}{t_{2}-t_{1}}\Bigr)
\cdots\\
\cdots
\exp\Bigl(-\frac{(x_{n}-x_{n-1})^{2}}{t_{n}-t_{n-1}}\Bigr)dx_{1}\cdots\,dx_{n-1},
\end{multline}
являющейся продолжением меры \eqref{eq:math:ex6} на всю нижнюю полуплоскость комплексного времени. Поскольку потенциальная энергия является непрерывной функцией координат и стремится к нулю на бесконечности, то подынтегральный функционал $$\displaystyle \exp\left(\int U\bigl(x(\tau)\bigr)\right)d\tau$$ ограничен и измерим \cite{bib:KPT}. Тогда континуальный интеграл является аналитической функцией комплексного времени области $-\pi<\varphi<0 $. В области $0<\varphi<\pi$ функция \eqref{eq:math:ex7} не ограничена для траекторий, уходящих на бесконечность. 
Чтобы обеспечить аналитичность интеграла в этой области, необходимо аналитически продолжить эту функцию через ось вещественного времени функцией
\begin{equation*}
\frac{1}{\bigl(\sqrt{2\pi\tau}\bigr)^{n}}\int_{\alpha_{1}}^{\beta_{1}}...
\int_{\alpha_{n}}^{\beta_{n}}\exp(-T_{1})\exp(-T_{2})
\cdots\exp(-T_{n})dx_{1}\cdots\,dx_{n-1}.
\end{equation*}
В этом выражении значение кинетической энергии $T$ равно $$\displaystyle \frac{m}{2}\Bigl(\frac{\Delta x}{\Delta \tau}\Bigr)^{2},$$ 
где $\Delta x$, $\Delta \tau$ "--- разница координат частицы и соответствующих им моментов времени на границах квазиинтервалов  \eqref{eq:math:ex5}. Таким образом, континуальный интеграл может быть определён как аналитическая функция для всех значений комплексного времени.

\smallskip

\Section[n]{Заключение} То обстоятельство, что квантовый интеграл по траекториям может быть определён для всей временной комплексной плоскости, позволяет привести его к вещественному виду не только для  $\varphi=-{\pi}/{2}$, но и для $\varphi={\pi}/{2}$. 
Первый случай соответствует обратной последовательности событий, тогда как второй "--- прямой. Наличие такой возможности позволяет исследовать изменение мер множеств траекторий при коллапсе (процессе, необратимом во времени).

\thanks{{\bf ORCID}

{\bf Наталья Витальевна Мелешко:} \url{http://orcid.org/0000-0001-8491-5539}

{\bf Aлексей Юрьевич Самарин:} \url{http://orcid.org/0000-0001-7640-3875}

}


\RusRef

\pagebreak

\begin{thebibliography}{10}


\Bibitem{bib:E}
\zmath{https://zbmath.org/?q=an:03109251}
\by von Neumann~J.
\book Mathematical foundations of quantum mechanics
\serial Investigations in Physics 
\vol 2
\publ Princeton
\publaddr Princeton Univ. Press
\totalpages xii+445 
\yr 1955

\Bibitem{bib:A} 
\paper Against `measurement'
\by Bell~J.
\jour Physics World
\issueinfo August
\yr 1990
\pages 33--40
%\elink \url{http://physicsworldarchive.iop.org/summary/pwa-xml/3/8/phwv3i8a26}
\moreref
\paper Against ``measurement''
\by Bell~J.~S.
\jour NATO ASI Series 
\vol 226
\yr 1990
\pages 17--31
\crossref{10.1007/978-1-4684-8771-8_3}


\Bibitem{bib:B} 
\paper ``Relative State'' Formulation of Quantum Mechanics
\by Everett~H.
\jour Rev. Mod. Phys.
\vol 29 
\issue 3 
\pages 454--462 
\yr 1957
\crossref{10.1103/revmodphys.29.454}


\RBibitem{bib:C} 
\by Менский~М.~Б.
\paper Квантовые измерения, феномен жизни и стрела времени: связи между ``тремя великими проблемами'' (по терминологии В.\,Л.~Гинзбурга)
\jour УФН
\yr 2007
\vol 177
\issue 4
\pages 415--425
\mathnet{http://mi.mathnet.ru/ufn458}
\crossref{10.3367/UFNr.0177.200704j.0415}
\adsnasa{http://adsabs.harvard.edu/cgi-bin/bib_query?2007PhyU...50..397M}
%\transl
%\jour Phys. Usp.
%\yr 2007
%\vol 50
%\issue 4
%\pages 397--407
%\crossref{http://dx.doi.org/10.1070/PU2007v050n04ABEH006241}
%\isi{http://gateway.isiknowledge.com/gateway/Gateway.cgi?GWVersion=2&SrcApp=PARTNER_APP&SrcAuth=LinksAMR&DestLinkType=FullRecord&DestApp=ALL_WOS&KeyUT=000248752700010}
%\scopus{http://www.scopus.com/record/display.url?origin=inward&eid=2-s2.0-34547851671}



\RBibitem{bib:D} 
\by Самарин~А.~Ю.
\paper Описание процесса перехода между состояниями дискретного спектра
\jour Вестн. Сам. гос. техн. ун-та. Сер. Физ.-мат. науки
\yr 2009
\issue 2(19)
\pages 226--230
\mathnet{http://mi.mathnet.ru/vsgtu721}
\crossref{10.14498/vsgtu721}


\RBibitem{bib:L}
\by Самарин~А.~Ю.
\paper Естественное пространство микрообъекта
\jour Вестн. Сам. гос. техн. ун-та. Сер. Физ.-мат. науки
\yr 2011
\issue 3(24)
\pages 117--128
\mathnet{http://mi.mathnet.ru/vsgtu911}
\crossref{10.14498/vsgtu911}

\Bibitem{bib:F} 
\paper Space-time approach to non-relativistic quantum mechanics
\by Feynman~R.~P.
\jour Rev. Mod. Phys.
\yr 1948
\vol 20
\issue 2
\pages 367--387
\crossref{10.1103/revmodphys.20.367}
\moreref
\paper Space-time approach to non-relativistic quantum mechanics
\by Feynman~R.~P.
\crossref{10.1142/9789812567635_0002}
\inbook  Feynman's Thesis — A New Approach to Quantum Theory
\yr 2005
\pages 71--109
\publ World Scientific Publ.
\publaddr Singapore

\Bibitem{bib:KV}
\by Feynman~R.~P., Hibbs~A.~R.
\book Quantum Mechanics and Path Integrals
\yr 1965
\publ McGraw-Hill
\publaddr New York
\totalpages 371+xii
\mathscinet{http://www.ams.org/mathscinet-getitem?mr=2797644}
\zmath{https://zbmath.org/?q=an:0176.54902}
\adsnasa{http://adsabs.harvard.edu/cgi-bin/bib_query?1965flp..book.....F}


\Bibitem{bib:V} 
\paper Quantum Particle Motion in Physical Space
\by Samarin~A.~Yu.
\jour  Advanced Studies in Theoretical Physics
\yr 2014
\vol 8
\issue 1
\pages 27--34
\crossref{10.12988/astp.2014.311136}
\arxiv \href{http://arxiv.org/abs/1407.3559}{1407.3559} [quant-ph]


\RBibitem{bib:PRL} 
\by Самарин~А.~Ю.
\paper Пространственная локализация квантовой частицы
\jour Вестн. Сам. гос. техн. ун-та. Сер. Физ.-мат. науки
\yr 2013
\issue 1(30)
\pages 387--397
\mathnet{http://mi.mathnet.ru/vsgtu1138}
\crossref{10.14498/vsgtu1138}

\RBibitem{bib:G}
\by Васильев~А.~Н.
\book Функциональные методы в квантовой теории поля и~статистики
\publaddr  Ленинград
\publ ЛГУ
\yr 1978
\totalpages 295

\RBibitem{bib:H}
\by Попов~В.~Н.
\book Континуальные интегралы в квантовой теории поля и~статистической физике
\publaddr  М.
\publ Атомиздат
\yr 1978
\totalpages 256

\Bibitem{bib:K}
\by Faddeev~L.~D., Slavnov~A.~A.
\book Gauge fields: introduction to quantum theory
\zmath{https://zbmath.org/?q=an:0486.53052}
\serial Frontiers in Physics
\vol 50
\publaddr  Reading, Mass.
\publ Benjamin/Cummings, Advanced Book Program
\totalpages xiii+232 
\yr 1980

\Bibitem{bib:EDK}
\by Zinn Justin~J.
\book Path Integrals in Quantum Mechanics
\yr 2004
\publ Oxford University Press
\publaddr Oxford
\totalpages 320+xiv
\crossref{10.1093/acprof:oso/9780198566748.001.0001}
\mathscinet{http://www.ams.org/mathscinet-getitem?mr=2724801}


\Bibitem{bib:KSV} 
\preprint  Macroscopic Body Motion in Terms of Quantum Evolution
\yr 2004
\by Samarin~A.~Yu.
\totalpages 5
\arxiv \href{http://arxiv.org/abs/1408.0340}{1408.0340} [quant-ph]


\Bibitem{bib:KPT}
\by Kac~M.
\book Probability and related topics in physical sciences
\serial Lectures in Applied Mathematics
\vol I
\publ Interscience Publ.
\publaddr London, New York 
\yr 1959 
\totalpages xiii+266 


\end{thebibliography}


\RuDate


\newpage



\makeentitle



\thanks{{\bf ORCID}

{\bf Natalia V. Meleshko:} \url{http://orcid.org/0000-0001-8491-5539}

{\bf Alexey Yu. Samarin:} \url{http://orcid.org/0000-0001-7640-3875}

}


\EngRef

\pagebreak


\begin{thebibliography}{10}


\Bibitem{bib:E:eng}
\zmath{https://zbmath.org/?q=an:03109251}
\by von Neumann~J.
\book Mathematical foundations of quantum mechanics
\serial Investigations in Physics 
\vol 2
\publ Princeton
\publaddr Princeton Univ. Press
\totalpages xii+445 
\yr 1955

\Bibitem{bib:A:eng} 
\paper Against `measurement'
\by Bell~J.
\jour Physics World
\issueinfo August
\yr 1990
\pages 33--40
%\elink \url{http://physicsworldarchive.iop.org/summary/pwa-xml/3/8/phwv3i8a26}
\moreref
\paper Against ``measurement''
\by Bell~J.~S.
\jour NATO ASI Series 
\vol 226
\yr 1990
\pages 17--31
\crossref{10.1007/978-1-4684-8771-8_3}


\Bibitem{bib:B:eng} 
\paper ``Relative State'' Formulation of Quantum Mechanics
\by Everett~H.
\jour Rev. Mod. Phys.
\vol 29 
\issue 3 
\pages 454--462 
\yr 1957
\crossref{10.1103/revmodphys.29.454}


\Bibitem{bib:C:eng} 
\paper Quantum measurements, the phenomenon of life, and time arrow: three great problems of physics (in Ginzburg's terminology) and their interrelation
\by Mensky~M.~B.
\jour Phys. Usp.
\yr 2007
\vol 50
\issue 4
\pages 397--407
\crossref{10.1070/PU2007v050n04ABEH006241}
\isi{http://gateway.isiknowledge.com/gateway/Gateway.cgi?GWVersion=2&SrcApp=PARTNER_APP&SrcAuth=LinksAMR&DestLinkType=FullRecord&DestApp=ALL_WOS&KeyUT=000248752700010}
\scopus{http://www.scopus.com/record/display.url?origin=inward&eid=2-s2.0-34547851671}



\Bibitem{bib:D:eng} 
\lang In Russian
\by Samarin~A.~Yu.
\paper Description of discrete spectrum states transition processes
\jour Vestn. Samar. Gos. Tekhn. Univ. Ser. Fiz.-Mat. Nauki
\yr 2009
\issue 2(19)
\pages 226--230
\crossref{10.14498/vsgtu721}


\Bibitem{bib:L:eng}
\lang In Russian
\by Samarin~A.~Yu.
\paper Natural space of the micro-object
\jour Vestn. Samar. Gos. Tekhn. Univ. Ser. Fiz.-Mat. Nauki
\yr 2011
\issue 3(24)
\pages 117--128
\crossref{10.14498/vsgtu911}

\Bibitem{bib:F:eng} 
\paper Space-time approach to non-relativistic quantum mechanics
\by Feynman~R.~P.
\jour Rev. Mod. Phys.
\yr 1948
\vol 20
\issue 2
\pages 367--387
\crossref{10.1103/revmodphys.20.367}
\moreref
\paper Space-time approach to non-relativistic quantum mechanics
\by Feynman~R.~P.
\crossref{10.1142/9789812567635_0002}
\inbook  Feynman's Thesis — A New Approach to Quantum Theory
\yr 2005
\pages 71--109
\publ World Scientific Publ.
\publaddr Singapore

\Bibitem{bib:KV:eng}
\by Feynman~R.~P., Hibbs~A.~R.
\book Quantum Mechanics and Path Integrals
\yr 1965
\publ McGraw-Hill
\publaddr New York
\totalpages 371+xii
\mathscinet{http://www.ams.org/mathscinet-getitem?mr=2797644}
\zmath{https://zbmath.org/?q=an:0176.54902}
\adsnasa{http://adsabs.harvard.edu/cgi-bin/bib_query?1965flp..book.....F}


\Bibitem{bib:V:eng} 
\paper Quantum Particle Motion in Physical Space
\by Samarin~A.~Yu.
\jour  Advanced Studies in Theoretical Physics
\yr 2014
\vol 8
\issue 1
\pages 27--34
\crossref{10.12988/astp.2014.311136}
\arxiv \href{http://arxiv.org/abs/1407.3559}{1407.3559} [quant-ph]


\Bibitem{bib:PRL:eng} 
\lang In Russian
\by Samarin~A.~Yu.
\paper Space localization of the quantum particle
\jour Vestn. Samar. Gos. Tekhn. Univ. Ser. Fiz.-Mat. Nauki
\yr 2013
\issue 1(30)
\pages 387--397
\crossref{10.14498/vsgtu1138}

\Bibitem{bib:G:eng}
\by Vasiliev~A.~N.
\book Functional Methods in Quantum Field Theory and Statistical Physics 
\yr 1998
\publ  Gordon and Breach Science Publ.
\publaddr Amsterdam
\totalpages  xiv+312 
\zmath{https://zbmath.org/?q=an:01446851}

\Bibitem{bib:H:eng}
\lang In Russian
\by Popov~V.~N.
\book Kontinual'nye integraly v kvantovoi teorii polia i~statisticheskoi fizike \rm [Path Integrals in Quantum Field Theory and Statistical Physics]
\publaddr  Moscow
\publ Atomizdat
\yr 1978
\totalpages 256

\Bibitem{bib:K:eng}
\by Faddeev~L.~D., Slavnov~A.~A.
\book Gauge fields: introduction to quantum theory
\zmath{https://zbmath.org/?q=an:0486.53052}
\serial Frontiers in Physics
\vol 50
\publaddr  Reading, Mass.
\publ Benjamin/Cummings, Advanced Book Program
\totalpages xiii+232 
\yr 1980

\Bibitem{bib:EDK:eng}
\by Zinn Justin~J.
\book Path Integrals in Quantum Mechanics
\yr 2004
\publ Oxford University Press
\publaddr Oxford
\totalpages 320+xiv
\crossref{10.1093/acprof:oso/9780198566748.001.0001}
\mathscinet{http://www.ams.org/mathscinet-getitem?mr=2724801}


\Bibitem{bib:KSV:eng} 
\preprint  Macroscopic Body Motion in Terms of Quantum Evolution
\yr 2004
\by Samarin~A.~Yu.
\totalpages 5
\arxiv \href{http://arxiv.org/abs/1408.0340}{1408.0340} [quant-ph]


\Bibitem{bib:KPT:eng}
\by Kac~M.
\book Probability and related topics in physical sciences
\serial Lectures in Applied Mathematics
\vol I
\publ Interscience Publ.
\publaddr London, New York 
\yr 1959 
\totalpages xiii+266 


\end{thebibliography}


\EngDate

\renewcommand{\baselinestretch}{0.9}

 \newpage
\Russian
%
% \tableofengcontents

%\end{document}
%
